
Machine learning reproducibility is affected by dataset evolution. When classifiers trained on the original ImageNet were tested on ImageNet-v2, accuracy decreased by 11-14 percentage points \citep{Rechtetal2019}. This performance degradation occurred despite no changes to model architecture or training procedures, indicating that distribution shifts between dataset versions can substantially affect model performance.

Existing dataset versioning systems—Data Version Control (DVC) and HuggingFace Datasets—track changes using opaque identifiers. A researcher examining version identifier \texttt{rev=a3f7b2d} cannot determine whether it represents a change requiring model retraining or a minor documentation update. Software engineering addresses this problem through semantic versioning (MAJOR.MINOR.PATCH), where version numbers indicate compatibility: MAJOR versions indicate breaking changes, MINOR versions indicate backward-compatible additions, and PATCH versions indicate backward-compatible fixes.

This work tested whether statistical drift detection could enable automated semantic versioning for datasets. The hypothesis was that fixed thresholds (7\% for MAJOR, 2\% for MINOR, 0.5\% for PATCH) applied to drift scores computed from Kolmogorov-Smirnov tests and Maximum Mean Discrepancy could classify version changes with ≥85\% accuracy, ≥70\% precision, and ≥85\% recall.

\textbf{Results:} The classifier achieved 28.57\% accuracy on 14 dataset pairs, with 25\% precision and 25\% recall for MAJOR change detection. These metrics fall below the random baseline (33\% for three classes) and substantially below the predefined success thresholds. The classifier missed 3 out of 4 MAJOR changes while incorrectly classifying 6 non-MAJOR changes as MAJOR.

\textbf{Root cause:} Drift scores ranged from 0.042 to 0.79 across datasets, an 18.8× variance. Multiple PATCH-labeled datasets exceeded the 0.07 MAJOR threshold, while some MAJOR-labeled datasets fell below it. This indicates that drift magnitude is dataset-relative: a score of 0.79 may represent normal variation for one dataset while 0.087 represents a significant shift for another.

\textbf{Contributions:}

\begin{enumerate}
\item Empirical falsification of fixed-threshold semantic versioning across 14 dataset pairs spanning vision and NLP domains.
\end{enumerate}

\begin{enumerate}
\item Quantitative measurement of cross-dataset drift variance (18.8× range), demonstrating that universal thresholds cannot accommodate dataset-specific baselines.
\end{enumerate}

\begin{enumerate}
\item Documentation of systematic failure modes: the classifier exhibited 75\% false negative rate for MAJOR changes and produced six false MAJOR predictions for one true positive.
\end{enumerate}

\begin{enumerate}
\item Identification of dataset coverage limitations affecting generalization claims, with only 1 vision dataset tested and 6 datasets unavailable due to API or access constraints.
\end{enumerate}

This negative result provides evidence that statistical drift magnitude alone, without dataset-specific calibration, is insufficient for semantic version classification. The paper presents detailed failure analysis and discusses implications for dataset versioning research.
